\scanpage{164}
\mspage{163}
VII.18

LEMMA VII.15. If $S$ is a semi-local domain and $S_{\uline M}$ is
a unique factorization domain for each maximal ideal $\uline M$ of $S$,
then $S$ is a unique factorization domain.

\uline{Proof:} Let $\unsure{\mathfrak p|\mathcal P}$ be a prime ideal in $S$. We will show that $\mathfrak p$ is
a principal ideal. Since $S$ is an integrally closed domain, this is
equivalent to saying that $S$ is a unique factorization domain. If $\uline M$
is a maximal ideal in $S$, $\mathfrak p_{\uline M}$ is either $S_{\uline M}$ or a prime ideal in $S_{\uline M}$.
Since $S_{\uline M}$ is presumed to be a unique factorization domain $\mathfrak p_{\uline M}$ is
a principal ideal (in either case). Thus $\mathfrak p_{\uline M}$ is a free $S_{\uline M}$ module
for each $\uline M$, and $\mathfrak p$ is therefore projective when considered an
$S$-module. Then the Remark in [4, p. 377] proves that $\mathfrak p$ is a
principal ideal in $S$.

LEMMA VII.16. Every regular semi-local domain is a unique
factorization domain.

\uline{Proof:} If $S$ is a regular semi-local domain, then $S_{\uline M}$ is a
regular local domain for each maximal ideal $\uline M$ of $S$. Auslander
and Buchsbaum have shown in [3] that every regular local domain is a
unique factorization domain.

THEOREM VII.17. If $A$ is a regular local domain, then
$\mathcal B(A)\subset\mathcal B(K)$.
